\documentclass[11pt]{article}

\usepackage[T1]{fontenc}
\usepackage[utf8]{inputenc}
\usepackage{amsmath}
\usepackage{amssymb}
\usepackage{amsthm}
\usepackage{geometry}
\geometry{a4paper, margin=25mm}

\newtheorem{assumption}{Assumption}
\newtheorem{remark}{Remark}

\newcommand{\E}{\mathbb{E}}
\newcommand{\Var}{\operatorname{Var}}
\newcommand{\dto}{\xrightarrow{d}}

\title{A Gaussian Maximum-Likelihood Estimator for $\gamma$\\
\large from per-scale sample means}
\author{\texttt{loglog\_validation} project notes}
\date{}

\begin{document}
\maketitle

\begin{abstract}
\noindent
This note derives a maximum-likelihood estimator (MLE) for the critical
exponent $\gamma$ that is \emph{not} the log-log OLS estimator already
implemented in \texttt{tools/loglog.py}. That estimator works by taking
$\log \overline{Y}_i$ as (approximately) Gaussian and fitting a line. Here we
instead take $\overline{Y}_i$ itself as (approximately) Gaussian, motivated
directly by the classical CLT applied to a sample mean, and derive the MLE for
$\gamma$ under that model. The two estimators are different objects, both
CLT-motivated, and this note also records exactly how they differ
(Remark~\ref{rmk:vs-ols}). This is a derivation to be checked, not yet
implemented.
\end{abstract}

\section{Setup and assumptions}

We are given data at $K$ distinct scales $i_1 < i_2 < \cdots < i_K$. At scale
$i_k$, we have $n_k$ i.i.d.\ samples $Y_k^{(1)}, \dots, Y_k^{(n_k)}$ of the
observable $Y_{i_k}$, and we write their sample mean as
\[
  \overline{Y}_k := \frac{1}{n_k}\sum_{\ell=1}^{n_k} Y_k^{(\ell)}.
\]
Samples at different scales are drawn independently (the project's
independence rule), so $\overline{Y}_1, \dots, \overline{Y}_K$ are independent
random variables.

\begin{assumption}[Relative variance, constant across scales]
\label{ass:relvar}
Write $\mu_k := \E\, Y_{i_k}$ for the (unknown) population mean at scale
$i_k$. We assume
\[
  \Var(Y_{i_k}) = \sigma^2 \mu_k^2
\]
for a single unknown constant $\sigma^2 > 0$, the same at every scale used.
This is the constant-variance regime of the article's Assumption~6
($\sigma_k^2 \to \sigma_\infty^2$; here we take $\sigma_k^2 \equiv \sigma^2$
exactly, consistent with $\xi_k := Y_{i_k}/\mu_k$ having $\Var(\xi_k) =
\sigma^2$ for every $k$ used).
\end{assumption}

\begin{assumption}[Leading-order power law]
\label{ass:powerlaw}
\[
  \mu_k = a_0\, i_k^{\gamma}
\]
for unknown constants $a_0 > 0$, $\gamma \in \mathbb{R}$ (the article's
eq.~232 with no finite-size correction term; see
Remark~\ref{rmk:corrections} for the extension).
\end{assumption}

\begin{assumption}[Gaussian approximation of the sample mean]
\label{ass:clt}
By the classical CLT (i.i.d., finite variance),
\[
  \sqrt{n_k}\,(\overline{Y}_k - \mu_k) \dto \mathcal{N}(0, \Var(Y_{i_k}))
  = \mathcal{N}(0, \sigma^2 \mu_k^2) \qquad \text{as } n_k \to \infty,
\]
and for the finite (but large) $n_k$ actually used we treat this as an
approximate equality in distribution:
\[
  \overline{Y}_k \ \approx\ \mathcal{N}\!\left(\mu_k,\ \frac{\sigma^2 \mu_k^2}{n_k}\right).
\]
\end{assumption}

The unknowns are $(\gamma, a_0, \sigma^2)$; the data are
$(i_k, n_k, \overline{Y}_k)_{k=1}^K$, all observed/known.

\section{The likelihood}

Under Assumption~\ref{ass:clt} and independence across $k$, the (approximate)
joint density of $(\overline{Y}_1,\dots,\overline{Y}_K)$ is
\[
  L(\gamma, a_0, \sigma^2)
  = \prod_{k=1}^K
    \left(2\pi\,\frac{\sigma^2\mu_k^2}{n_k}\right)^{-1/2}
    \exp\!\left[-\,\frac{n_k(\overline{Y}_k - \mu_k)^2}{2\sigma^2 \mu_k^2}\right],
  \qquad \mu_k = a_0 i_k^\gamma.
\]
Taking logarithms,
\begin{equation}
\label{eq:loglik}
  \ell(\gamma, a_0, \sigma^2)
  = -\frac{K}{2}\log(2\pi) - \frac{K}{2}\log\sigma^2
    - \sum_{k=1}^K \log \mu_k + \frac{1}{2}\sum_{k=1}^K \log n_k
    - \frac{1}{2\sigma^2}\sum_{k=1}^K \frac{n_k(\overline{Y}_k-\mu_k)^2}{\mu_k^2}.
\end{equation}

\section{Score equations}

We maximize \eqref{eq:loglik} over $(\gamma, a_0, \sigma^2)$ by setting its
partial derivatives to zero.

\subsection{With respect to $\sigma^2$}

\[
  \frac{\partial \ell}{\partial \sigma^2}
  = -\frac{K}{2\sigma^2} + \frac{1}{2\sigma^4}\sum_{k=1}^K \frac{n_k(\overline{Y}_k-\mu_k)^2}{\mu_k^2}
  = 0
\]
gives, in closed form,
\begin{equation}
\label{eq:sigmahat}
  \hat\sigma^2(\gamma, a_0) = \frac{1}{K}\sum_{k=1}^K \frac{n_k(\overline{Y}_k-\mu_k)^2}{\mu_k^2}.
\end{equation}

\subsection{With respect to $\mu_k$, then $\gamma$ and $a_0$}

Write $r_k := (\overline{Y}_k - \mu_k)/\mu_k$ for the relative residual at
scale $k$ (so $\overline{Y}_k = \mu_k(1+r_k)$). Differentiating the $k$-th
term of \eqref{eq:loglik} with respect to $\mu_k$ (holding $\sigma^2$ fixed),
using the quotient rule on $(\overline{Y}_k-\mu_k)^2/\mu_k^2$:
\[
  \frac{\partial}{\partial \mu_k}
  \left[\frac{(\overline{Y}_k-\mu_k)^2}{\mu_k^2}\right]
  = \frac{-2(\overline{Y}_k-\mu_k)\mu_k^2 - (\overline{Y}_k-\mu_k)^2 \cdot 2\mu_k}{\mu_k^4}
  = -\frac{2\overline{Y}_k(\overline{Y}_k-\mu_k)}{\mu_k^3},
\]
where the last equality uses $\mu_k + (\overline{Y}_k - \mu_k) = \overline{Y}_k$. Hence
\[
  \frac{\partial \ell}{\partial \mu_k}
  = -\frac{1}{\mu_k}
    + \frac{n_k\,\overline{Y}_k(\overline{Y}_k-\mu_k)}{\sigma^2 \mu_k^3}
  = \frac{1}{\mu_k}\left[\frac{n_k(1+r_k)r_k}{\sigma^2} - 1\right],
\]
using $\overline{Y}_k(\overline{Y}_k-\mu_k)/\mu_k^2 = (1+r_k)r_k$.

Since $\mu_k = a_0 i_k^\gamma$, we have
\[
  \frac{\partial \mu_k}{\partial \gamma} = \mu_k \log i_k,
  \qquad
  \frac{\partial \mu_k}{\partial a_0} = \frac{\mu_k}{a_0}.
\]
By the chain rule,
\begin{align}
  \frac{\partial \ell}{\partial \gamma}
  &= \sum_{k=1}^K \frac{\partial \ell}{\partial \mu_k}\cdot \mu_k \log i_k
   = \sum_{k=1}^K \log i_k
     \left[\frac{n_k(1+r_k)r_k}{\sigma^2} - 1\right] = 0,
   \label{eq:score-gamma}\\[4pt]
  \frac{\partial \ell}{\partial a_0}
  &= \frac{1}{a_0}\sum_{k=1}^K
     \left[\frac{n_k(1+r_k)r_k}{\sigma^2} - 1\right] = 0
   \ \Longleftrightarrow\
   \sum_{k=1}^K \left[\frac{n_k(1+r_k)r_k}{\sigma^2} - 1\right] = 0.
   \label{eq:score-a0}
\end{align}

\begin{remark}[A clean form for the $a_0$ equation]
Write $w_k := n_k r_k^2 = n_k(\overline{Y}_k-\mu_k)^2/\mu_k^2$, so that
\eqref{eq:sigmahat} reads $\sum_k w_k = K\hat\sigma^2$. Expanding
$n_k(1+r_k)r_k = n_k r_k + w_k$, equation~\eqref{eq:score-a0} becomes
$\sum_k (n_k r_k + w_k)/\sigma^2 = K$. At the joint stationary point
$\sigma^2=\hat\sigma^2$, so $\sum_k w_k/\hat\sigma^2 = K$ and this reduces to
\begin{equation}
\label{eq:score-a0-clean}
  \sum_{k=1}^K n_k\,\frac{\overline{Y}_k - \mu_k}{\mu_k} = 0,
\end{equation}
a weighted normal equation with no $\sigma^2$ dependence left --- useful as an
exact identity to check any numerical solution against. The $\gamma$
equation~\eqref{eq:score-gamma} does \emph{not} simplify the same way (the
extra $\log i_k$ weight blocks the same cancellation), so it must be kept in
its $\sigma^2$-dependent form, or solved jointly with
\eqref{eq:sigmahat}--\eqref{eq:score-a0-clean}.
\end{remark}

\section{Reduced form for implementation: the profile log-likelihood}

Equations \eqref{eq:sigmahat}, \eqref{eq:score-gamma}, \eqref{eq:score-a0} are
three equations in three unknowns with no closed-form solution for
$(\gamma, a_0)$ in general. The standard and most practical way to solve them
numerically is to \emph{profile out} $\sigma^2$: substitute
$\hat\sigma^2(\gamma,a_0)$ from \eqref{eq:sigmahat} back into
\eqref{eq:loglik}, which by construction removes $\sigma^2$ as a free
variable and leaves an objective in $(\gamma, a_0)$ alone.

Substituting \eqref{eq:sigmahat} into the last term of \eqref{eq:loglik}:
\[
  \frac{1}{2\hat\sigma^2}\sum_{k=1}^K \frac{n_k(\overline{Y}_k-\mu_k)^2}{\mu_k^2}
  = \frac{1}{2\hat\sigma^2}\cdot K\hat\sigma^2 = \frac{K}{2},
\]
so
\[
  \ell\big(\gamma, a_0, \hat\sigma^2(\gamma,a_0)\big)
  = -\frac{K}{2}\log(2\pi) - \frac{K}{2}\log\hat\sigma^2(\gamma,a_0)
    - \sum_{k=1}^K \log \mu_k + \frac{1}{2}\sum_{k=1}^K \log n_k - \frac{K}{2}.
\]
Dropping the additive constants that do not depend on $(\gamma, a_0)$,
maximizing the profile log-likelihood is equivalent to \emph{minimizing}
\begin{equation}
\label{eq:objective}
  \boxed{\
  \Phi(\gamma, a_0) \;:=\; K \log \hat\sigma^2(\gamma, a_0) \;+\; 2\sum_{k=1}^K \log \mu_k(\gamma, a_0),
  \qquad \mu_k = a_0 i_k^\gamma,\ \ \hat\sigma^2 \text{ as in \eqref{eq:sigmahat}}.
  \ }
\end{equation}
This is a plain function of two real variables $(\gamma, a_0)$ and is what
should actually be handed to a numerical optimizer (e.g.\
\texttt{scipy.optimize.minimize}), starting from the existing OLS-on-log
estimate (\texttt{tools/loglog.py}'s \texttt{gamma\_all\_points}) as an
initial guess for $\gamma$, and e.g.\ $\overline{Y}_{k_{\min}}/i_{k_{\min}}^{\hat\gamma_0}$
as an initial guess for $a_0$.

\begin{remark}[Consistency check]
By the envelope theorem, $\partial \Phi/\partial\gamma = 0$ and
$\partial\Phi/\partial a_0 = 0$ are exactly equations \eqref{eq:score-gamma}
and \eqref{eq:score-a0} evaluated at $\sigma^2 = \hat\sigma^2(\gamma,a_0)$ ---
minimizing \eqref{eq:objective} and solving the original score equations are
the same problem. A numerical minimizer's solution should therefore also
satisfy the clean identity~\eqref{eq:score-a0-clean}; checking that is a
cheap correctness test on any implementation.
\end{remark}

\begin{remark}[Extension to the $J$-order expansion]
\label{rmk:corrections}
If instead $\mu_k = a_0 i_k^\gamma \exp\!\big(\sum_j a_j i_k^{-\omega_j}\big)$
(the article's full eq.~232, correction exponents $\omega_j$ treated as
known), every derivative above goes through unchanged: the correction factor
does not depend on $\gamma$ or $a_0$, so $\partial \mu_k/\partial\gamma =
\mu_k \log i_k$ and $\partial\mu_k/\partial a_0 = \mu_k/a_0$ still hold
exactly. Equations \eqref{eq:score-gamma}--\eqref{eq:objective} are therefore
unchanged in form; only the definition of $\mu_k$ changes.
\end{remark}

\begin{remark}[How this differs from the OLS-on-log estimator]
\label{rmk:vs-ols}
\texttt{tools/loglog.py}'s existing estimators work by treating
$\log\overline{Y}_k$ as approximately Gaussian around $\log\mu_k$ (a
delta-method consequence of Assumption~\ref{ass:clt}, valid when the relative
fluctuation $\sigma/\sqrt{n_k}$ is small) and fitting an \emph{unweighted}
OLS line to $(\log i_k, \log \overline{Y}_k)$. The estimator derived here
instead applies the CLT directly to $\overline{Y}_k$, without a log
transform, and consequently weighs each scale by $n_k/\mu_k^2$ (visible in
\eqref{eq:score-a0-clean} and in $\hat\sigma^2$) rather than uniformly. The
two coincide asymptotically as $n_k \to \infty$ (delta method), but are
different estimators at finite $n_k$, and comparing them empirically was the
point of building \texttt{compare\_methods} in the first place.
\end{remark}

\section{Summary for implementation}

Given $(i_k, n_k, \overline{Y}_k)_{k=1}^K$:
\begin{enumerate}
\item Compute $\mu_k(\gamma,a_0) = a_0 i_k^\gamma$ and $\hat\sigma^2(\gamma,a_0)$
      via \eqref{eq:sigmahat}.
\item Minimize $\Phi(\gamma,a_0)$ of \eqref{eq:objective} numerically over
      $(\gamma, a_0)$, e.g.\ with \texttt{scipy.optimize.minimize}, starting
      from the OLS-on-log estimate.
\item Report $(\hat\gamma, \hat a_0)$ at the minimizer, and
      $\hat\sigma^2 = \hat\sigma^2(\hat\gamma,\hat a_0)$ from
      \eqref{eq:sigmahat}.
\item Sanity check: \eqref{eq:score-a0-clean},
      $\sum_k n_k(\overline{Y}_k-\hat\mu_k)/\hat\mu_k \approx 0$, should hold
      at the reported solution to floating-point/optimizer tolerance.
\end{enumerate}

\end{document}
